\documentclass[a4paper,11pt]{article}
\usepackage[utf8]{inputenc}
\usepackage[T1]{fontenc}
\usepackage[polish]{babel}
\usepackage{geometry}
\usepackage{hyperref}
\usepackage{listings}
\usepackage{xcolor}
\usepackage{booktabs}
\usepackage{pgfplots}
\pgfplotsset{compat=1.18}
\geometry{margin=2.5cm}
\lstset{
  basicstyle=\ttfamily\footnotesize,
  backgroundcolor=\color{gray!5},
  frame=single,
  breaklines=true,
  keywordstyle=\color{blue},
  commentstyle=\color{gray},
  showstringspaces=false,
  numbers=left,
  numberstyle=\tiny,
  xleftmargin=2em,
  framexleftmargin=1.5em,
}
\begin{document}
\title{Raport optymalizacji systemu rezerwacji salonu}
\author{System rezerwacji \texttt{FastAPI + MySQL}}
\date{16 czerwca 2026}
\maketitle
\section{Cel}
Celem raportu jest przedstawienie kolejnych etapów optymalizacji zapytania i architektury systemu rezerwacji fryzjerskiej na podstawie wykonanych testów wydajnościowych. W raporcie opisano trzy etapy: usunięcie zliczania rekordów, analizę planu zapytania z EXPLAIN i wymuszenie kolejności JOIN, oraz rozdzielenie odczytów i zapisów na dwie bazy MySQL.

\section{Środowisko testowe}
System został zbudowany jako aplikacja FastAPI z bazą MySQL uruchamianą w Dockerze. W projekcie znajdują się następujące elementy:
\begin{itemize}
  \item \texttt{main.py} \textendash FastAPI + logika endpointów,
  \item \texttt{db.py} \textendash zarządzanie pulami połączeń MySQL,
  \item \texttt{docker-compose.yml} \textendash usługi \texttt{db\_primary}, \texttt{db\_replica} i \texttt{app}.
\end{itemize}
Do testów użyto JMetera i czterech iteracji obciążenia danymi z próbami od 10 do 250 użytkowników.

\section{Etap 1: usunięcie zliczania rekordów}
Pierwsza poprawka polegała na usunięciu dodatkowego zapytania, które wykorzystywało \texttt{COUNT(*)}. To zliczanie było niepotrzebne dla paginacji i powodowało dodatkowy koszt na dużym zbiorze danych.

W kodzie miał miejsce fragment:
\begin{lstlisting}[language=python]
offset = (page - 1) * limit
\end{lstlisting}

Zamiast dodatkowego \texttt{COUNT(*)} pobieraliśmy dane bezpośrednio z \texttt{LIMIT} i \texttt{OFFSET}.

Dzięki temu zredukowano liczbę operacji agregacji i poprawiono wydajność odczytu. Ta zmiana znalazła się w zestawieniu \texttt{summary\_po\_zmianie.csv}.

\section{Etap 2: analiza EXPLAIN i wymuszenie kolejności JOIN}
Kolejny etap to analiza planu wykonania zapytania przy użyciu instrukcji \texttt{EXPLAIN SELECT}. Plan wykazał niekorzystny przebieg zapytania, gdzie MySQL skanował tabelę \texttt{salon} zamiast wykorzystywać indeksy po stronie tabeli \texttt{time\_slot}.

Aby wymusić lepszy plan, zastosowano \lstinline{STRAIGHT_JOIN}, który każe silnikowi MySQL łączyć tabele w zadanej kolejności. W \texttt{main.py} zmieniono konstrukcje JOIN ze standardowych na:
\begin{lstlisting}[language=SQL]
FROM time_slot ts
STRAIGHT_JOIN salon ON salon.salon_id = ts.salon_id
STRAIGHT_JOIN hairdresser ON hairdresser.hairdresser_id = ts.hairdresser_id
\end{lstlisting}

Po tej zmianie pojawił się wynik z \texttt{summary\_po\_explain.csv}, gdzie widoczna jest przyrostowa poprawa w próbach 100 i wyższych, przy zachowaniu niskiego odsetka błędów.

\section{Etap 3: rozdzielenie baz na primary/replika}
Ostatni etap to przebudowa infrastruktury na dwie bazy MySQL: \texttt{db\_primary} dla zapisów i \texttt{db\_replica} dla odczytów.
Konfiguracja w \texttt{docker-compose.yml} obejmuje:
\begin{itemize}
  \item \texttt{db\_primary}: port 3307, binlog, server-id=1,
  \item \texttt{db\_replica}: port 3308, read-only, relay log, server-id=2,
  \item \texttt{app}: zmienne \texttt{DB\_WRITE\_HOST=db\_primary} i \texttt{DB\_READ\_HOST=db\_replica}.
\end{itemize}

Kod read/write split w \texttt{db.py}:
\begin{lstlisting}[language=Python]
DB_WRITE_HOST = os.getenv("DB_WRITE_HOST", "127.0.0.1")
DB_WRITE_PORT = int(os.getenv("DB_WRITE_PORT", 3307))
DB_READ_HOST = os.getenv("DB_READ_HOST", "127.0.0.1")
DB_READ_PORT = int(os.getenv("DB_READ_PORT", 3308))
\end{lstlisting}

Dzięki takiej separacji odczyty mogły być obsługiwane przez replikę bez obciążania bazy zapisu.

\section{Wyniki}
Poniżej przedstawione są średnie czasy odpowiedzi dla kolejnych serii testów.

\begin{table}[ht]
\centering
\caption{Porównanie czasów \texttt{Average} dla eksperymentów}
\label{tab:wyniki}
\begin{tabular}{@{}rrrrr@{}}
\toprule
Label & Przed zmianą & Po zmianie & Po EXPLAIN & Dwie bazy \\
\midrule
10 & 11 & 9 & 10 & 9 \\
25 & 10 & 8 & 9 & 7 \\
50 & 13 & 9 & 9 & 8 \\
75 & 24 & 13 & 13 & 9 \\
100 & 110 & 77 & 42 & 15 \\
125 & 181 & 122 & 114 & 38 \\
150 & 260 & 214 & 195 & 111 \\
175 & 342 & 287 & 270 & 168 \\
200 & 446 & 366 & 357 & 227 \\
225 & 527 & 424 & 428 & 311 \\
250 & 626 & 501 & 470 & 355 \\
\midrule
TOTAL & 361 & 290 & 276 & 186 \\
\bottomrule
\end{tabular}
\end{table}

\begin{figure}[ht]
\centering
\begin{tikzpicture}
\begin{axis}[
    width=\textwidth,
    height=9cm,
    grid=major,
    xlabel={Liczba próbek (Label)},
    ylabel={Średni czas odpowiedzi [ms]},
    legend pos=north west,
    ymin=0,
    ymax=700,
    xtick={10,25,50,75,100,125,150,175,200,225,250},
    x tick label style={rotate=45,anchor=east},
]
\addplot+[mark=*,blue] table [col sep=comma, x=Label, y=Average] {
Label,Average
10,11
25,10
50,13
75,24
100,110
125,181
150,260
175,342
200,446
225,527
250,626
};
\addlegendentry{Przed zmianą}
\addplot+[mark=triangle*,red] table [col sep=comma, x=Label, y=Average] {
Label,Average
10,9
25,8
50,9
75,13
100,77
125,122
150,214
175,287
200,366
225,424
250,501
};
\addlegendentry{Po usunięciu count}
\addplot+[mark=square*,orange] table [col sep=comma, x=Label, y=Average] {
Label,Average
10,10
25,9
50,9
75,13
100,42
125,114
150,195
175,270
200,357
225,428
250,470
};
\addlegendentry{Po EXPLAIN / STRAIGHT\_JOIN}
\addplot+[mark=pentagon*,teal] table [col sep=comma, x=Label, y=Average] {
Label,Average
10,9
25,7
50,8
75,9
100,15
125,38
150,111
175,168
200,227
225,311
250,355
};
\addlegendentry{Dwie bazy}
\end{axis}
\end{tikzpicture}
\caption{Porównanie średnich czasów odpowiedzi dla czterech etapów optymalizacji.}
\end{figure}

\begin{figure}[ht]
\centering
\begin{tikzpicture}
\begin{axis}[
    width=100mm,
    height=100mm,
    grid=major,
    xlabel={Liczba próbek (Label)},
    ylabel={Średni czas odpowiedzi [ms]},
    legend pos=north west,
    ymin=0,
    ymax=30,
    xmin=10,
    xmax=100,
    xtick={10,25,50,75,100},
    x tick label style={rotate=45,anchor=east},
]
\addplot+[mark=*,blue] table [col sep=comma, x=Label, y=Average] {
Label,Average
10,11
25,10
50,13
75,24
100,110
};
\addlegendentry{Przed zmianą}
\addplot+[mark=triangle*,red] table [col sep=comma, x=Label, y=Average] {
Label,Average
10,9
25,8
50,9
75,13
100,77
};
\addlegendentry{Po usunięciu count}
\addplot+[mark=square*,orange] table [col sep=comma, x=Label, y=Average] {
Label,Average
10,10
25,9
50,9
75,13
100,42
};
\addlegendentry{Po EXPLAIN / STRAIGHT\_JOIN}
\addplot+[mark=pentagon*,teal] table [col sep=comma, x=Label, y=Average] {
Label,Average
10,9
25,7
50,8
75,9
100,15
};
\addlegendentry{Dwie bazy}
\end{axis}
\end{tikzpicture}
\caption{Wykres 100×100 dla wartości do 100, pokazujący realne czasy przed załamaniem krzywej.}
\end{figure}

\section{Wnioski}
Testy pokazują, że największą poprawę przyniosła architektura z dwoma bazami oraz separacja odczytów i zapisów. Zmiana zliczania wpłynęła na umiarkowaną redukcję opóźnienia, a analiza EXPLAIN z wymuszeniem kolejności JOIN dała dalsze zauważalne przyspieszenie.

W szczególności:
\begin{itemize}
  \item usunięcie dodatkowego zapytania \texttt{COUNT} obniżyło średni czas z 361 ms do 290 ms,
  \item wymuszenie kolejności JOIN przy użyciu \texttt{STRAIGHT\_JOIN} obniżyło średni czas do 276 ms,
  \item finalna konfiguracja z dwiema bazami obniżyła średni czas do 186 ms i poprawiła przepustowość odczytów.
\end{itemize}

Optymalizacja zapytań i architektury bazy danych w tej kolejności okazała się skuteczna: najpierw zmniejszono koszty pojedynczego zapytania, a następnie poprawiono skalowalność na poziomie infrastruktury.

\end{document}
